% body-de.tex -- shared content of the three package specimens (design-packages-A, 2026-10-09).
% Semantic macros only; each pkgN.tex defines them.  Depth devices follow layered-docs/REPORT.md 5.4
% (Kern / Vertiefung / Fundament; run-in label, dependency statement, return line).  Physics checked by hand:
%   rho = (1 + r.sigma)/2, E = (1 + n.sigma)/2 with |n| = 1  =>  tr(rho E) = (1 + r.n)/2;
%   E_n = - mu e^4 / (2 (4 pi eps0)^2 hbar^2 n^2) ~ -13,6 eV / n^2;  E_3 - E_2 = 1,89 eV.
\pkgchapter{6}{Die Bornsche Regel}

\pkgkicker{Zoomstufe 1 \textendash\ ein einzelnes Quantensystem}

Eine physikalische Theorie liefert Daten für drei Dinge: Präparationen, Dynamik und Messungen.
Eine Präparation nennen wir \enquote{Zustand}; im Quantenfall beschreiben wir sie durch einen
Dichteoperator $\rho\geq 0$ mit $\operatorname{tr}\rho=1$.\pkgsidenote{Für ein Qubit ist jeder
Dichteoperator von der Form $\rho=\tfrac12(\mathbb{1}+\pkgvec{r}\cdot\pkgvec{\sigma})$ mit
$\lvert\pkgvec{r}\rvert\leq 1$. Reine Zustände liegen auf dem Rand, $\lvert\pkgvec{r}\rvert=1$.}
Ein Detektor, der \enquote{klick} macht oder schweigt, wird durch einen Effekt $E$ mit
$0\leq E\leq\mathbb{1}$ beschrieben. Die Theorie sagt nur Wahrscheinlichkeiten voraus, und zwar
genau eine Zahl für jedes Paar aus Zustand und Effekt.

\begin{pkgpostulate}{P4}{Bornsche Regel}
Wird ein System im Zustand $\rho$ präpariert und mit dem Effekt $E$ getestet, so spricht der
Detektor mit der Wahrscheinlichkeit
\begin{equation}
  p(E\,\vert\,\rho)=\operatorname{tr}(\rho E)
\end{equation}
an. Insbesondere gilt $0\leq p\leq 1$, und der triviale Effekt $E=\mathbb{1}$ spricht immer an.
\end{pkgpostulate}

Für ein Qubit lässt sich die Regel zeichnen (Abb.~\ref{fig:born}a). Mit
$\rho=\tfrac12{(\mathbb{1}+\pkgvec{r}\cdot\pkgvec{\sigma})}$ und dem Projektor
$E=\tfrac12{(\mathbb{1}+\pkgvec{n}\cdot\pkgvec{\sigma})}$, $\lvert\pkgvec{n}\rvert=1$, folgt
\begin{equation}
  \operatorname{tr}(\rho E)=\tfrac12\bigl(1+\pkgvec{r}\cdot\pkgvec{n}\bigr),
  \qquad
  E_n=-\frac{\mu e^4}{2(4\pi\varepsilon_0)^2\hbar^2}\,\frac{1}{n^2}\approx-\frac{\qty{13.6}{\electronvolt}}{n^2}.
\end{equation}
Die Wahrscheinlichkeit ist also die Projektion des Bloch-Vektors auf die Richtung des
Detektors, verschoben und halbiert. Die Zeitentwicklung $\rho\mapsto U_t\,\rho\,U_t^{*}$ mit
$U_t=\mathrm{e}^{-\mathrm{i}Ht/\hbar}$ dreht den Vektor; sie ändert $\lvert\pkgvec{r}\rvert$ nicht.

\begin{figure}[h]
  % fig-born.tex -- shared specimen figure.  Styles (state, effect, energy, env, ladder1..4, lbl,
% forbidden, axisstyle, rulegrey) are defined by each pkgN.tex from its token table.
% Geometry: r = 0.75 at 50 deg, n = e_x;  r.n = 0.75 cos 50 = 0.482;  p = (1 + 0.482)/2 = 0.741.
% Hydrogen: E_n = -13.6/n^2 eV -> -13.6, -3.40, -1.51, -0.85;  0.5 cm per eV for n >= 2, n = 1 drawn
% below an explicit axis break (its true position would be 6.8 cm below zero).
\begin{center}
\begin{tikzpicture}[x=1cm,y=1cm]
  % ---------------- (a) Bloch disc
  \begin{scope}[scale=1.3]
    \fill[forbidden, even odd rule] (0,0) circle (1.16) (0,0) circle (1.0);
    \draw[axisstyle] (-1.25,0) -- (1.32,0);
    \draw[axisstyle] (0,-1.22) -- (0,1.3);
    \draw[boundary] (0,0) circle (1.0);
    \draw[effect, -{Stealth[length=5pt]}] (0,0) -- (1.0,0) node[lbl, below right, inner sep=1pt] {$\pkgvec{n}$};
    \draw[state, -{Stealth[length=5pt]}] (0,0) -- (50:0.75) node[lbl, above right, inner sep=1pt] {$\pkgvec{r}$};
    \fill[statefill] (50:0.75) circle (0.035);
    \draw[proj] (50:0.75) -- (0.482,0);
    \draw[effect, line width=2.2pt] (0,0) -- (0.482,0);
    \node[lbl, below, inner sep=2pt] at (0.24,-0.02) {$\pkgvec{r}\cdot\pkgvec{n}$};
    \node[lbltext, anchor=south west] at (-1.18,1.16) {(a)};
    \node[lbltext, anchor=north west, align=left] at (1.0,-0.80) {\pkgsmall reine\\[-1pt]\pkgsmall Zustände};
    \draw[leader] (1.0,-0.86) -- (-40:1.0);
  \end{scope}
  % ---------------- (b) hydrogen ladder, 0.5 cm per eV above the break; n = 1 below a marked break
  \begin{scope}[xshift=5.2cm, yshift=1.35cm, yscale=0.92]
    \node[lbltext, anchor=south west] at (-1.5,0.25) {(b)};
    \draw[axisstyle] (0,-2.05) -- (0,0.45) node[lbl, left] {$E$};
    \draw[axisline] (0,-2.75) -- (0,-3.75);
    \draw[axisline] (-0.12,-2.13) -- (0.12,-2.03) (-0.12,-2.27) -- (0.12,-2.17);   % axis break
    \draw[rulegrey] (-0.08,0) -- (0.08,0) node[lbl, left=3pt] {$0$};
    \foreach \e/\c in {-3.40/ladder2, -1.511/ladder3, -0.85/ladder4}
      \draw[\c, line width=1.3pt] (0.25,{0.5*\e}) -- (2.05,{0.5*\e});
    \draw[ladder1, line width=1.3pt] (0.25,-3.6) -- (2.05,-3.6);
    \node[lbl, right] at (2.1,-3.6) {$n=1$};
    \node[lbl, right] at (2.1,{0.5*-3.40}) {$n=2$};
    \node[lbl, right] at (2.1,{0.5*-1.511}) {$n=3$};
    \node[lbl, right] at (2.1,{0.5*-0.85}) {$n=4$};
    \draw[energy, -{Stealth[length=5pt]}] (0.75,{0.5*-1.511}) -- (0.75,{0.5*-3.40});
    \node[lbl, right, inner sep=2pt] at (0.8,{0.5*-2.45}) {$\qty{1.89}{\electronvolt}$};
    \node[lbl, left] at (-0.05,{0.5*-3.40}) {$\qty{-3.40}{\electronvolt}$};
    \node[lbl, left] at (-0.05,-3.6) {$\qty{-13.6}{\electronvolt}$};
  \end{scope}
\end{tikzpicture}
\end{center}

  \caption{\pkgcaptionlead{Die Wahrscheinlichkeit ist eine Projektion.} (a) Zustand $\pkgvec{r}$ im
  Bloch-Kreis (Schnitt durch die Bloch-Kugel), Detektorrichtung $\pkgvec{n}$; die Projektion
  liefert $p=\tfrac12(1+\pkgvec{r}\cdot\pkgvec{n})\approx 0{,}74$. Außerhalb des Kreises liegen
  keine Zustände (schraffiert). (b) Wasserstoff: die Niveaus $E_n$ für $n=1,\dots,4$; der Übergang
  $3\to2$ setzt $\qty{1.89}{\electronvolt}$ frei.}
  \label{fig:born}
\end{figure}

\begin{pkgtracktwo}{6.V2}{Zustände als Funktionale}
\pkgdepends{Voraussetzung: Kern 6.1. Gebraucht erst bei unendlich vielen Freiheitsgraden.}
Allgemein ist ein Zustand ein positives, normiertes lineares Funktional $\omega$ auf der Algebra
$\mathcal{B}(\mathcal{H})$, also $\omega(A^{*}A)\geq 0$ und $\omega(\mathbb{1})=1$. Die
\emph{normalen} Zustände sind genau die Funktionale $\omega(A)=\operatorname{tr}(\rho A)$ mit einem
Dichteoperator $\rho$ der Spurklasse. Ist $\dim\mathcal{H}=\infty$, so gibt es auch singuläre Zustände.
\pkgreturn{Im Folgenden genügen Dichteoperatoren.}
\end{pkgtracktwo}

Messgeräte sind also Elemente des Dualraums: ein Effekt ordnet jedem Zustand eine Zahl zwischen
$0$ und $1$ zu, linear in $\rho$. Das ist der Grund, warum Mischungen von Präparationen zu
Mischungen von Wahrscheinlichkeiten führen.
